\scanpage{29}
\mspage{23}

Les ``pré quintuplets de Quillen clos'', associés
à $W$, correspondent au choix de
parties closes à gauche $TC$ et $C$,
s'insérant dans le diagramme (I) comme
montré \inserted{en $\mathrm I'$} ci-dessous, et donnant lieu
par transposition à $F$ et $TF$ comme dans
$\mathrm{II}'$ ci-dessous

\diagram{\begin{tikzcd}[ampersand replacement=\&,row sep=small,column sep=small]
 \& \& W^{\mathrm{univ}}\arrow[ddrr,hook] \& \& \\
 \& TC\arrow[ur,hook]\arrow[ddrr,hook] \& \& \& \\
 (\mathrm I')\qquad(\mathrm{Fib}_W)^{\unsure{v|*}}\arrow[ur,hook]\arrow[ddrr,hook] \& \& \& \& \mathrm{Cof}_W \\
 \& \& \& C\arrow[ur,hook] \& \\
 \& \& (W_{\mathrm{univ}})^{\unsure{*|k}}\arrow[ur,hook,no head] \& \&
\end{tikzcd}}

\diagram{\begin{tikzcd}[ampersand replacement=\&,row sep=small,column sep=small]
 \& \& (W^{\mathrm{univ}})_*\arrow[dl,hook] \& \& \\
 \& F\arrow[dl,hook] \& \& \& \\
 (\mathrm{II}')\qquad\mathrm{Fib}_W \& \& \& \& (\mathrm{Cof}_W)_{\unsure{\alpha|*}}\arrow[uull,hook]\arrow[dl,hook,thick] \\
 \& \& \& TF\arrow[uull,hook]\arrow[dl,hook,thick] \& \\
 \& \& W_{\mathrm{univ}}\arrow[uull,hook] \& \&
\end{tikzcd}}\unsure{;}

le ``\unsure{hic}'' étant donc de choisir $TC$ et $C$
assez grandes pour qu'on ait
\[
 \begin{aligned}
 (TC)_*&\subset\mathrm{Fib}_W,&(\deleted{\unsure{TF}}C)_*&\subset W_{\mathrm{univ}}\quad-\\
 \parallel &&\parallel\\
 F&&TF
 \end{aligned}
\]
% archive sheet number: 28
